% Original synthetic authoring demonstration. Not a scientific submission.
% Compile from the skill root after rendering reports/gallery, if TeX is installed:
% xelatex -output-directory=reports assets/templates/manuscript-en.tex
% Requires Times New Roman. XeLaTeX/Tectonic; do not use pdfLaTeX.
% Check final layout, references, journal requirements, and length separately.
\documentclass[aps,prl,reprint]{revtex4-2}
\usepackage{amsmath,amssymb,bm,graphicx}
\usepackage{mathspec}
\setmainfont{Times New Roman}
\setsansfont{Times New Roman}
\setmathsfont(Digits,Latin){Times New Roman}
\setmathrm{Times New Roman}
\setmathsf{Times New Roman}
\setmathtt{Times New Roman}
\graphicspath{{reports/gallery/}{../../reports/gallery/}}
\begin{document}
\title{Two analytical response curves: a synthetic authoring demonstration}
\author{Example Author}
\affiliation{Synthetic demonstration; no institutional affiliation is asserted}
\begin{abstract}
Two analytical response curves illustrate a figure workflow. Their differences
are defined by their generating functions and do not establish an empirical mechanism.
\end{abstract}
\maketitle
\textit{Construction.—}
We compare two analytical response curves to illustrate fixed-size figure export.
Figure~\ref{fig:trend} shows their values with illustrative deviations.
Both curves are specified by construction; no experimental measurement or
mechanistic test is implied.
\begin{figure}[tb]
\centering
\includegraphics[width=\columnwidth]{01-trend.png}
\caption{Two synthetic response series. Circles and squares identify the series,
and solid and dashed curves show their generating functions. Bars show fixed
illustrative deviations of 0.025 arbitrary units, not measured statistical uncertainty.}
\label{fig:trend}
\end{figure}
% Add only verified references when adapting this demonstration to actual research.
% End Matter is separate from the core. Follow the current PRL instructions for it;
% never move evidence essential to the central argument out of the core merely to fit.
\end{document}
